\pdfbookmark[0]{Streszczenie}{streszczenie.1}
%\phantomsection
%\addcontentsline{toc}{chapter}{Streszczenie}
%%% Poniższe zostało niewykorzystane (tj. zrezygnowano z utworzenia nienumerowanego rozdziału na abstrakt)
%%%\begingroup
%%%\setlength\beforechapskip{48pt} % z jakiegoś powodu była maleńka różnica w położeniu nagłówka rozdziału numerowanego i nienumerowanego
%%%\chapter*{\centering Abstrakt}
%%%\endgroup
%%%\label{sec:abstrakt}
%%%Lorem ipsum dolor sit amet eleifend et, congue arcu. Morbi tellus sit amet, massa. Vivamus est id risus. Sed sit amet, libero. Aenean ac ipsum. Mauris vel lectus. 
%%%
%%%Nam id nulla a adipiscing tortor, dictum ut, lobortis urna. Donec non dui. Cras tempus orci ipsum, molestie quis, lacinia varius nunc, rhoncus purus, consectetuer congue risus. 
%\mbox{}\vspace{2cm} % można przesunąć, w zależności od długości streszczenia
\begin{abstract}
Celem pracy jest stworzenie projektu oraz implementacja aplikacji ułatwiającej współpracę przy grupowych projektach studenckich oraz pozyskiwanie potencjalnych inwestorów.
Dziedzina problemu została podzielona na kilka podproblemów takich jak hosting repozytoriów kodu źródłowego, tworzenie sprintów i zadań,
proces rekrutacji do projektów oraz możliwość odkrywania projektów przez biznes. Na podstawie podziału problemu zdefiniowano wymagania funkcjonalne i~niefunkcjonalne oraz przypadki użycia.
Projekt zakłada stworzenie dwóch aplikacji, serwerowej i klienckiej webowej.
Aplikację serwerową zaimplementowano zgodnie z elementami podejścia projektowego \textit{DDD} \cite{evans2015domain} z wykorzystaniem wzorca projektowego \textit{CQRS} \cite{CQRS}.
Aplikacja kliencka została oparta o \textit{React.js} w oparciu o specyfikację \textit{Material Design} w wersji trzeciej \cite{MaterialDesign}. Największym wyzwaniem projektu była obsługa repozytoriów \textit{Git},
która wymgała od Autora głębokiej analizy dokumentacji \textit{Git} oraz opracowania odpowiedniej struktury systemu przechowywania ich. 
Do obu aplikacji utworzono testy jednostkowe, integracyjne i akceptacyjne w formule \textit{E2E}. Testowano między innymi zachowanie repozytorium po wypchnięciu na serwer
oraz mechanizmy uwierzytelniania. Testy zaimplementowano za pomocą narzędzi \textit{TUnit} oraz \textit{Playwright}.
Implementację projektu zrealizowano zgodnie z powstałymi wymaganiami, wyciągnięto wnioski na temat repozytoriów oraz zdefiniowano przyszłe kierunki rozwoju,
takie jak na przykład wyszukiwarka projektów, tagi i~implementacja wewnętrznego mechanizmu komunikacji między biznesem a wykonawcami projektu. 
\end{abstract}
\mykeywords

{
\selectlanguage{english}
\begin{abstract}
The main goal of this thesis is creation of project and implementation of application which will faciliate
collaboration on student group projects and and attract potential investors. Field of the problem was divided into smaller
subproblems consisting of repository hosting, sprint and task creation, recrutation process and project exploration for business. Based on subproblem definition
unctional requirements, non-functional requirements and use cases were defined. The project assumes that there will be created client web application and server one.
The server app was created with elements of Domain-Driven-Design and heavy usage of \textit{CQRS} \cite{CQRS}. Client app was created with
\textit{React.js}, largely based on \textit{Material Design} specification in third version. The biggest challenge of implemetation was to handle \textit{Git} repositories appropriately,
which demanded from Author to deeply analyze \textit{Git} documentation and work out appropriate solution to the system of repository storage. Unit, integration and acceptance in E2E formula tests were created for both applications.
Testing consisted of for example behaviour of repository after pushing it to the server and authorization mechanism. The tests were implemented 
using \textit{TUnit} and \textit{Playwright} frameworks. Implementation of the project was made according to previously defined requirements, the conclusions were made and possible future path was defined.
Future path can possibly consist of project search engine, tags and implementation of internal communication system between business and project members.
\end{abstract}
\mykeywords
}
